\section{Koneksi Database: Pengenalan MySQL dan MySQLi}

\textbf{MySQL} adalah sistem manajemen basis data relasional (\textit{RDBMS}) open-source yang paling banyak digunakan bersama PHP. MySQL menggunakan bahasa SQL (\textit{Structured Query Language}) untuk operasi data dan mendukung fitur seperti \textit{stored procedures}, \textit{triggers}, \textit{views}, dan \textit{transactions}. Kombinasi PHP dan MySQL (sering disebut sebagai bagian dari stack \textbf{LAMP} --- Linux, Apache, MySQL, PHP) telah menjadi fondasi ribuan aplikasi web, dari blog sederhana hingga platform e-commerce berskala besar \cite{mysqlDocs}.

PHP menyediakan ekstensi \textbf{MySQLi} (MySQL Improved) sebagai pengganti ekstensi \code{mysql\_*} yang sudah dihapus sejak PHP 7.0. MySQLi menawarkan dua gaya pemrograman: \textbf{prosedural} (mirip fungsi lama) dan \textbf{OOP} (menggunakan class dan object). Gaya OOP lebih direkomendasikan karena lebih bersih, konsisten, dan sejalan dengan paradigma OOP yang dipelajari di bab ini. MySQLi mendukung \textit{prepared statements}, \textit{multiple statements}, dan \textit{transactions} --- fitur yang tidak ada di ekstensi lama \cite{phpManual}.

\begin{figure}[ht]
\centering
\begin{tikzpicture}[
  box/.style={draw, rounded corners, minimum width=3.2cm, minimum height=1.2cm, align=center, font=\small, fill=#1},
  arrow/.style={-Stealth, thick}
]
  \node[box=blue!15] (php) {Aplikasi PHP\\{\footnotesize (server-side)}};
  \node[box=orange!20, right=2cm of php] (mysqli) {Ekstensi\\MySQLi};
  \node[box=green!15, right=2cm of mysqli] (mysql) {MySQL\\Server};
  \node[box=yellow!20, below=1.5cm of mysql] (db) {Database\\{\footnotesize (tabel, data)}};

  \draw[arrow] (php) -- node[above, font=\footnotesize]{\$conn = new mysqli()} (mysqli);
  \draw[arrow] (mysqli) -- node[above, font=\footnotesize]{TCP/Socket} (mysql);
  \draw[arrow] (mysql) -- (db);
  \draw[arrow] (db) -- (mysql);
  \draw[arrow] (mysql) -- node[below, font=\footnotesize]{Result set} (mysqli);
  \draw[arrow] (mysqli) -- node[below, font=\footnotesize]{\$result} (php);
\end{tikzpicture}
\caption{Arsitektur koneksi PHP -- MySQLi -- MySQL}
\end{figure}

\begin{webcode}[language=PHP, caption={Koneksi MySQLi gaya OOP dengan error handling}]
<?php
// Konfigurasi - idealnya di file terpisah (config.php)
$host     = 'localhost';
$username = 'root';
$password = '';
$database = 'belajar_php';

// Membuat koneksi (OOP)
$conn = new mysqli($host, $username, $password, $database);

// Cek koneksi
if ($conn->connect_error) {
  die("Koneksi gagal: " . $conn->connect_error);
}

// Set charset UTF-8 (penting untuk karakter Indonesia)
$conn->set_charset("utf8mb4");

echo "Koneksi ke database '{$database}' berhasil!";

// ... lakukan query ...

// Tutup koneksi setelah selesai
$conn->close();
?>
\end{webcode}

Beberapa praktik terbaik dalam mengelola koneksi database: (1) simpan kredensial database di file konfigurasi terpisah yang tidak dapat diakses dari web (misalnya di luar \code{public\_html/}); (2) selalu set charset ke \code{utf8mb4} untuk mendukung karakter Unicode penuh; (3) cek \code{connect\_error} setelah koneksi dan berikan pesan error yang informatif; (4) tutup koneksi dengan \code{\$conn->close()} setelah selesai (meskipun PHP menutupnya otomatis di akhir skrip); (5) di lingkungan produksi, gunakan \textit{connection pooling} atau \textit{persistent connections} untuk performa yang lebih baik. Jangan pernah menampilkan detail koneksi database (host, username) ke pengguna jika koneksi gagal \cite{phpRightWay}.

